一份网关日志摊在屏幕上，几十万行，每行有用的信息只有一个时间戳。你想从中读出点别的：现在的容量够不够、这次超时告警是不是真的异常、如果流量翻倍队列会堆到多长。这些问题的形状是一样的——事件零散地落在一条时间轴上，而你要对下一段时间说点什么。

计数过程给这件事一套记法。记 \(N(t)\) 为到时刻 \(t\) 为止累计发生的事件数，路径从零出发、只往上跳、跳幅是整数。整条路径由事件时刻完全决定，于是同一份数据可以从两头读：圈定一段时间数里面有几个事件，或者盯住某个事件序号问它什么时候到。这两种读法在本章会来回切换。

主干只有一句话。若到达满足三条极简假设——不相交时段互不影响（独立增量）、等长时段的统计行为相同（平稳增量）、同一瞬间不会挤进两个事件——过程就被唯一确定了，速率 \(\lambda\) 是仅有的自由参数。此时计数服从 Poisson、相邻间隔服从 Exponential、第 \(k\) 次到达服从 Gamma。这三件事不是三条各自成立的结论，它们是同一个过程被三种方式读出来的样子；分布层面的对照见\hyperref[sec:05-Probability/06-Common-Distributions/03]{``Poisson 计数、Exponential 等待与 Gamma 时间''}，本章关心的是过程层面。

\begin{center}
\begin{tikzpicture}[x=1cm,y=1cm,font=\footnotesize,>=stealth]
  \draw[black!45,fill=blue!6] (0,0) rectangle (3.4,2.3);
  \node at (1.7,1.95) {三条假设};
  \node at (1.7,1.35) {独立增量};
  \node at (1.7,0.85) {平稳增量};
  \node at (1.7,0.35) {瞬时不重叠};
  \draw[->,black!65] (3.55,1.15) -- (4.35,1.15);
  \draw[black!45,fill=blue!14] (4.5,0.55) rectangle (7.2,1.75);
  \node at (5.85,1.35) {Poisson 过程};
  \node at (5.85,0.85) {唯一参数 \(\lambda\)};
  \draw[->,black!65] (7.35,1.15) -- (8.15,1.15);
  \draw[black!45,fill=blue!6] (8.3,0) rectangle (11.7,2.3);
  \node at (10.0,1.95) {三种读法};
  \node at (10.0,1.35) {计数 \(N(t)\)：Poisson};
  \node at (10.0,0.85) {间隔 \(W_r\)：Exponential};
  \node at (10.0,0.35) {到达 \(T_r\)：Gamma};
  \draw[->,black!55,dashed] (1.0,-0.15) -- (1.0,-1.05);
  \node[right,black!70] at (1.15,-0.6) {放松平稳增量 \(\Rightarrow\) 非齐次 Poisson};
  \draw[->,black!55,dashed] (1.0,-1.25) -- (1.0,-2.15);
  \node[right,black!70] at (1.15,-1.7) {放松单位跳跃 \(\Rightarrow\) 复合 Poisson};
  \draw[->,black!55,dashed] (8.6,-0.15) -- (8.6,-1.05);
  \node[right,black!70] at (8.75,-0.6) {放松指数间隔 \(\Rightarrow\) 更新过程};
\end{tikzpicture}

\emph{本章的骨架。左边三条假设合起来只留下一个参数；右边是同一个过程的三种读数。要动脑的地方在虚线上：每放松一条假设，模型往哪个方向走。}
\end{center}

先在离散时隙里把两种读法建立起来，再让时隙长度趋于零得到连续时间的版本，然后逐条松绑。速率随时间变化（白天与凌晨的流量差一个数量级）松的是平稳增量；一次到达携带一个随机大小（一笔订单的金额、一个请求的字节数）松的是单位跳跃；间隔带记忆（硬盘随役龄老化、用户等到不耐烦就走）松的是指数间隔本身，这时要换成更新过程。

\subsection{本章篇目}

\begin{itemize}
\tightlist
\item \hyperref[sec:06-Random-Processes/02-Counting-Processes/01]{``Bernoulli 过程与离散到达''}：在按秒轮询的粗粒度世界里，先看清``数事件''与``量时间''是同一件事的两面；
\item \hyperref[sec:06-Random-Processes/02-Counting-Processes/02]{``Poisson 过程与独立增量''}：三条假设如何锁死整个过程，以及非齐次与复合两个最常用的推广；
\item \hyperref[sec:06-Random-Processes/02-Counting-Processes/03]{``到达时刻、指数等待与条件均匀性''}：把无记忆性讲透——它为什么让排队模型变得可算，以及怎样发现它在你的数据上不成立；
\item \hyperref[sec:06-Random-Processes/02-Counting-Processes/04]{``更新过程与长期更新率''}：间隔一般化之后还剩什么，以及随机时刻观察为什么会系统性地偏向长周期。
\end{itemize}

四篇按依赖排列，前两篇是主线。若你手上已经有一份到达数据、只想判断 Poisson 是否够用，可以直接跳到第三篇的危险率诊断，再按结论回补。本章的模型在\hyperref[ch:050]{``连续时间 Markov 链与排队：不仅要知道跳到哪，还要知道何时跳''}里会被当作零件反复调用，那里的状态空间之所以能压得那么小，靠的正是这里的无记忆性。
